\section[Named switch centers and the fixed-P4 physical consumer]{Named switch centers and the fixed-$P_4$ physical consumer}
\label{sec:switch-consumer}
This section integrates Notes22--23 \cite{ansatz22,ansatz23} and
Assessments16--17. Its compact input is Theorem~\ref{thm:compact-positive}
on $[63,65]$, with $J=2$ and $P=P_4$; no eventual late-time onset
is evaluated at $64$. The parent \cite{parent} supplies the exact
machinery in its positive-evolution section:
\begin{itemize}
\setlength{\itemsep}{0pt}
\item Lemma \texttt{lem:holomorphic-weak-density};
\item the current identity \texttt{eq:positive-weak-current};
\item the signed stress identity \texttt{eq:positive-signed-stress};
\item the balance identity \texttt{eq:positive-general-residual-balance};
\item the weak correction \texttt{eq:positive-weak-residual-correction};
\item Proposition \texttt{prop:positive-residual-equations}.
\end{itemize}
The derivations below consume these identities at the corrected
model's actual collars and pay all terms. This is a production
integration of same-model assessed source results, not independent
specialist validation.

\begin{proposition}[Named-center sensitivity and finite-height shifts]
\label{prop:named-switch-centres}
For fixed $\ell,h,\eta>0$, on $63\le t\le65$ at sufficiently large
$X$, preserve the exact historical center
$b_{\rm lin}(t_0,X)=0$, where
$b_{\rm lin}=\log(\epsilon_g/|G(t,s_X)|)$. With
$\lambda=\log X$ and $w_X=128/(\pi X\lambda)$,
\[
 B_0=b_{{\rm lin},t}(t_0,X)
 =-\frac{\pi X\lambda^2}{512}[1+O(1/\lambda)]
\]
holds uniformly on every fixed $Mw_X$ neighborhood with the derivative
and second-derivative bounds in \eqref{eq:ansatz-nc15}. Every named
one-arm modulus $A_0,A_1,A_J,A_S,A_4,A_{\rm exp}$ has a unique
comparison center there, and its paid finite-height shift is
\eqref{eq:ansatz-nc17}; all these displacements are $o(w_X)$.
Replacing a one-arm modulus by twice that modulus adds $\log2$ to
its logarithmic offset and gives the distinct shift
\eqref{eq:ansatz-nc19}. The inherited simplified-layer error remains
$O(1)$ absolute and $O(1/\log X)$ normalized.
\end{proposition}
The following argument gives the amplitude, direct-time, neighborhood
and layer payments for this proposition.

\subsection{Named one-arm amplitudes and the actual rounding neighborhood}
\label{sec:named-centres-proof}
\subsubsection{Fixed contract and actual compact input}
\label{subsec:nc-1}

Fix \(\ell,h,\eta>0\) and observation \(|y|\le 1.\) Put \(r=\ell+2, b=1+h+2,\) and use the full rectangles \(|\Re z-X|\le r, |\Im z|\le b\) with one further analytic unit. Throughout \(63\le t\le 65\) and X tends to positive infinity; evenness gives the opposite-height charts. Constants and finite onsets may depend on these fixed data. Set

\[
L=\log(X/(4\pi)),\quad \lambda=\log X,\quad \Lambda=1+\lambda,
\quad s_X=(1+iX)/2,
\]
\[
M(w)=w(w-1)\pi^{-w/2}\Gamma(w/2)/16,\quad\alpha=M'/M.
\tag{NC1}\label{eq:ansatz-nc1}
\]

The complete canonical \(H, H_{t}=-H_{zz}\) and H(0,z)=xi((1+iz)/2)/8 are unchanged. So are physical \(q=\epsilon_{g}^{2}\) and \(K=k_{\ell} k_{h},\) unnormalized density/current, additive potentials and actual covariance. We do not use an eventual large-t theorem at \(t=64.\) The input here is the actual compact theorem in Ansatz Note20, RB1--13, itself using the fixed-order theorem of Note19. Its P4 bound is \(C A \Lambda^{d} X^{-7}\) for spatial orders \(0\le d\le 3, C A \Lambda^{2} X^{-7}\) for direct first time, and \(C A X^{-7}\) for the actual heat residual. \(A=|S(s)|+|S(1-s)|\) is a positive arm envelope. Every full arithmetic tail, connector, middle and gamma end is included before absorption; \(63 \log2/4>7\) is the compact arithmetic gate.

Choose a single fixed \(V=\Sigma+3,\) where \(\Sigma\) bounds \(|\Re w|\) for \(w=s,1-s\) on the largest rectangle. In particular \(\Re(w+V)>2.\) Define

\[
G(t,w)=M(w)e^{t\alpha(w)^2/4},\qquad
J(t,w)=(\pi t)^{-1/2}\int_V^\infty e^{-v^2/t}M(w+v)\,dv,
\]
\[
p=w+(t/2)\alpha(p),\quad B=1-(t/2)\alpha'(p),\qquad
S(t,w)=M(p)e^{-(p-w)^2/t}B^{-1/2}.
\tag{NC2}\label{eq:ansatz-nc2}
\]

Use precisely the branch and gates of parent Note30 and compact RB1; the square root is near one. V is fixed in t and z. The reference J depends on this decomposition and is not an intrinsic amplitude of H. On the real axis define positive one-arm moduli

\[
A_{lin}=|G(t,s_X)|,\quad A_J=|J(t,s_X)|,\quad A_S=|S(t,s_X)|,
\]
\[
A_0=|C_+(t,X)|=A_o(t,X)/2,\quad
A_1=|C_+(t,X)(1+D_+(t,X))|.
\tag{NC3}\label{eq:ansatz-nc3}
\]

\(C_{\pm}\) and \(D_{\pm}\) are exactly RB6 of Note20, including unwrapped phase 7pi/8 and \(b_{\pm}=\Log(z/(4\pi))/2\) $\pm$ \(i \pi/4.\) For the fixed P4 model one may also take \(A_4=|S(t,s_{X})(1+C_{2}+C_{4})(t,s_{X})|,\) with the finite coefficients of Note19 FO8. All these arms are nonzero after an onset: their ratios to nonzero G or C tend to one. No logarithm of H or of a vanishing cosine is taken.

\subsubsection{Paid amplitude and direct time comparisons}
\label{subsec:nc-2}

From compact RB8/RB11, on the full larger holomorphic buffers,

\[
\log(J/G)=U+O(\Lambda^4/X^2),\quad
\log(S/G)=U+O(\Lambda^4/X^2),\qquad
U=t\alpha'/4+t^2\alpha^2\alpha'/8.
\tag{NC4}\label{eq:ansatz-nc4}
\]

The logarithms are the branches of the near-one ratios. The direct time remainder in either identity is \(O(\Lambda^{6}/X^{2}).\) This is established by differentiating the fixed Gaussian kernel for J, and the exact implicit identity for S, as in RB8a/RB11. It is not a derivative of an O allowance. To recall the payments: gamma cubic/quadratic-square error is \(C \Lambda^{4}/X^{2};\) arithmetic is \(C X^{-t \log2/4};\) original middle and missing lower Gaussian are \(C \exp(-2L^{2}+C(1+L));\) both infinite gamma ends are \(C \exp(-c X^{2/3}).\) Their direct kernel derivatives retain the same exponentials with at most two additional logarithmic moment powers.

The sectorial gamma inputs are CB4 of Note03. Their primary source, \href{https://dlmf.nist.gov/5.11}{DLMF5.11}, was rechecked on 26 October 2026 for its sector and logarithm conventions. With \(b_+=L/2+i \pi/4\) on the real axis,

\[
\log M(s_X)=\log C_+(0,X)+i/(24X)+O(X^{-2}),
\]
\[
\alpha(s_X)=b_+-7i/(2X)+O(X^{-2}),\quad
\alpha'(s_X)=-i/X+O(X^{-2}).
\tag{NC5}\label{eq:ansatz-nc5}
\]

Define \(g(t)=7\pi t/16\) and \(h(t,X)=\pi t^{2} L/32.\) Expanding exact holomorphic quantities before taking real parts gives

\[
\log(G/C_+)=i/(24X)-7itb_+/(4X)+O(\Lambda/X^2),
\]
\[
\Re U=h/X+O(\Lambda^2/X^2),\quad
\Re D_+=(g+h)/X.
\tag{NC6}\label{eq:ansatz-nc6}
\]

The determinant \(t \alpha'/4\) is retained in U; its leading value is imaginary on the real axis. Also \(\log|1+D|=\Re D+O(|D|^{2}),\) with \(|D|\le C \Lambda^{2}/X.\) Consequently \(\Delta_{A}=\log(A/A_{lin})\) satisfies:

\begingroup\small
\begin{longtable}{@{}>{\raggedright\arraybackslash}p{0.25\textwidth} >{\raggedright\arraybackslash}p{0.22\textwidth} >{\raggedright\arraybackslash}p{0.43\textwidth}@{}}
\toprule
Named modulus A & \(\Delta_{A}\) & Direct \(\partial_{t} \Delta_{A}\)\\
\midrule
\endhead
\(A_0,\) elementary one arm & $-g/X+O(\Lambda/X^2)$ & $-7\pi/(16X)+O(\Lambda/X^2)$\\
\(A_1,\) corrected elementary & $h/X+O(\Lambda^4/X^2)$ & $\pi tL/(16X)+O(\Lambda^6/X^2)$\\
\(A_{J},\) fixed truncated gamma arm & $h/X+O(\Lambda^4/X^2)$ & $\pi tL/(16X)+O(\Lambda^6/X^2)$\\
\(A_{S},\) full gamma saddle & $h/X+O(\Lambda^4/X^2)$ & $\pi tL/(16X)+O(\Lambda^6/X^2)$\\
\(A_4, P_{4}\) one arm & $h/X+O(\Lambda^4/X^2)$ & $\pi tL/(16X)+O(\Lambda^6/X^2)$\\
\bottomrule
\end{longtable}
\endgroup

For the last row, FO8 gives \(C_{2}=O(X^{-3}), C_{4}=O(X^{-5}),\) with direct-time derivatives of the same orders on this compact t-domain (use \(p_{t}=\alpha/(2B), |p|\) comparable to \(X, \alpha^{(j)}=O(X^{-j})).\) Thus it is within the displayed allowance. The real positive exponential rewrite \(A_{\exp}=(A_{o}/2)\exp((g+h)/X)\) has the same corrected row within its paid error; it is not identified exactly with \(A_1.\) Its holomorphic arm is \(C_+(t,z) \exp([g(t)+\pi t^{2} \Log(z/(4\pi))/32]/z),\) so its modulus on real X is \(A_{\exp}\) and the later phase comparison applies to it. Replacing any one-arm modulus A by 2A adds log2 to \(\Delta\) exactly. In particular \(A_{o}=2A_0,\) not \(A_0,\) and \(2A_1\) is the corresponding corrected cosine amplitude.

For clarity the full buffered field comparisons being consumed are:

\begingroup\small
\begin{longtable}{@{}>{\raggedright\arraybackslash}p{0.23\textwidth} >{\raggedright\arraybackslash}p{0.24\textwidth} >{\raggedright\arraybackslash}p{0.24\textwidth} >{\raggedright\arraybackslash}p{0.18\textwidth}@{}}
\toprule
Two-arm field F & Field error divided by comparable envelope \(E_{o}\) & Spatial jet order \(d\le 3\) & Direct first-time error\\
\midrule
\endhead
G(s)+G(1-s) & \(C \Lambda^{2}/X\) & \(C \Lambda^{2+d}/X\) & \(C \Lambda^{4}/X\)\\
J(s)+J(1-s) & \(C e_{J}\) & \(C \Lambda^{d} e_{J}\) & \(C \Lambda^{2} e_{J}\)\\
S(s)+S(1-s) & \(C X^{-3}\) & \(C \Lambda^{d} X^{-3}\) & \(C \Lambda^{2} X^{-3}\)\\
\(C_++C_-\) & \(C \Lambda^{2}/X\) & \(C \Lambda^{2+d}/X\) & \(C \Lambda^{4}/X\)\\
F1=sum \(C_{\pm}(1+D_{\pm})\) & \(C \Lambda^{4}/X^{2}\) & \(C \Lambda^{4+d}/X^{2}\) & \(C \Lambda^{6}/X^{2}\)\\
P4 & \(C X^{-7}\) & \(C \Lambda^{d} X^{-7}\) & \(C \Lambda^{2} X^{-7}\)\\
\bottomrule
\end{longtable}
\endgroup

Here
\[
 e_J=X^{-t\log2/4}+\exp[-2L^2+C(1+L)]+\exp(-cX^{2/3}).
\]
Constants can be enlarged on the fixed compact domain. Choose $c>0$
smaller once and enlarge the middle constant once to absorb the
polynomial weights from kernel differentiation. The $J$ comparison
retains complete ordinary arithmetic and the bounded middle. For $S$
use the compact $J=0$ allowance $X^{-3}$ from Theorem~\ref{thm:compact-positive}.
Cauchy disks of radius $c/\Lambda$ pay the displayed jets. The $G$
time comparison uses its explicit logarithm, $H_t=-H_{zz}$ and its
actual residual; $J$ uses the fixed kernel and $F_1,P_4$ their
separately proved direct-time bounds. Field error and residual remain distinct.

In particular the truncated arm has exact forcing

\[
(\partial_t-\tfrac14\partial_{ww})J
 =\frac{e^{-V^2/t}}{4\sqrt{\pi t}}
  [M'(w+V)+(2V/t)M(w+V)],
\tag{NC7}\label{eq:ansatz-nc7}
\]

with the superalgebraic lower-Gaussian/middle payment relative to G. For F1 use the exact RB10 residual; for P4 the exact FO30 identity at \(J=2\) and its \(C A X^{-7}\) allowance. None of these fields is declared to solve the homogeneous heat equation merely because it approximates H.

\subsubsection{Sensitivity on an actual rounding neighborhood}
\label{subsec:nc-3}

Preserve the historical definition, exactly,

\[
b_{lin}(t,X)=\log(\epsilon_g(t,X)/A_{lin}(t,X)),\qquad
b_{lin}(t_0,X)=0,\quad t_0=t_{c,lin}(X).
\tag{NC8}\label{eq:ansatz-nc8}
\]

The fixed regulator of parent Note20 FB1--2 is \(\epsilon_{b}=\eta C(1+t^{2})^{-1/4}, C=\Phi(0)\sqrt{\pi}/2,\) and, above64,

\[
\epsilon_g=E_0(t)e^{-\beta(t)X^2/2},\quad
E_0=\epsilon_b^{1-\chi}(M_0/\kappa)^\chi,
\quad\beta=\chi/\kappa.
\tag{NC9}\label{eq:ansatz-nc9}
\]

\(M_{0},\kappa\) are the original real source-saddle data of parent Note19; they are not the gamma factor M or its complex saddle p. For \(\delta=t-64\) in (0,1], the exact flat switch gives

\[
\log\beta=-64/\delta+f(\delta),\qquad
f(0)=1-\log\kappa(64),
\tag{NC10}\label{eq:ansatz-nc10}
\]

where f extends smoothly to0 with its first two derivatives bounded. Indeed \(\log \chi=-64/\delta+1/(1-\delta/64) -\log(1+\exp[-64/\delta+1/(1-\delta/64)]).\) The last term is flat at0; subtract \(\log \kappa.\) Thus directly, without differentiating a center allowance,

\[
\beta'/\beta=64/\delta^2+f',\quad
\beta''/\beta=(64/\delta^2+f')^2-128/\delta^3+f''.
\tag{NC11}\label{eq:ansatz-nc11}
\]

All required derivatives of log E0 are bounded on this fixed interval. The exact time derivative is

\[
b_{lin,t}=(\log E_0)'-\tfrac12 X^2\beta'
                          -\tfrac14\Re\alpha(s_X)^2.
\tag{NC12}\label{eq:ansatz-nc12}
\]

At64 b is \(\pi X/8+O(\Lambda^{2})>0,\) and at each fixed small \(\delta>0\) it is negative of order \(X^{2}.\) NC12 is strictly negative in a fixed sufficiently small right neighborhood of64 for large \(X: \beta'>0, \Re \alpha^{2}\) is positive of order \(\lambda^{2},\) and (log E0)' is bounded. This reproves local existence and uniqueness of t0. The equation at t0 and Stirling imply

\[
\beta(t_0)=\pi/(4X)+O(\Lambda^2/X^2),\qquad
t_0-64=\frac{64}{\lambda+c_0+O(1/\lambda)},\quad
c_0=\log\frac{4e}{\pi\kappa(64)}.
\tag{NC13}\label{eq:ansatz-nc13}
\]

The second formula locates the neighborhood; it is not substituted for the exact t0. Set

\[
w_X=128/(\pi X\lambda),\qquad
\mathcal N_M=\{|t-t_0|\le M w_X\},
\tag{NC14}\label{eq:ansatz-nc14}
\]

for each fixed \(M>0.\) Uniformly there, \(\delta\) is comparable to64/lambda. NC11 yields \(\beta(t)/\beta(t_{0})=1+O_{M}(\lambda/X),\) and hence

\[
B_0:=b_{lin,t}(t_0,X)
 =-\frac{\pi X\lambda^2}{512}[1+O(1/\lambda)],
\]
\[
b_{lin,t}(t,X)=B_0[1+O_M(\lambda/X)],\qquad
|b_{lin,tt}(t,X)|\le C_M X\Lambda^4.
\tag{NC15}\label{eq:ansatz-nc15}
\]

For the last bound \(\log A_{lin}\) is affine in t, so its second derivative is zero; NC11 bounds \(X^{2} \beta''.\) This supplies neighborhood control at the physical rounding scale, with the essential X factor. It does not differentiate JT30 or its O term.

\subsubsection{Centers, finite-height first corrections and error payments}
\label{subsec:nc-4}

For each named A define \(b_{A}=b_{lin}-\Delta_{A}\) and its comparison center \(t_{A}\) by \(b_{A}(t_{A},X)=0\) in \(N_{M}.\) Write \(m=c X \lambda^{2}\) for a uniform lower bound on \(-b_{{\rm lin},t}\) there, \(D=|\Delta_{A}(t_{0})|,\) and \(H_{A}=\sup_{NM} |\partial_t\Delta_A|.\) All one-arm rows have \(D+H_{A}\le C \Lambda/X\) after enlarging the onset; the same rows multiplied by a fixed positive constant have \(D=O(1).\) At the two endpoints of \(N_{M}, b_{lin}\) has opposite signs and magnitude comparable to \(M \lambda,\) whereas \(\Delta\) is bounded. Also \(b_{A,t}<0.\) Therefore each center exists uniquely there, for large enough X.

The mean-value bound \(|t_{A}-t_{0}|\le C D/m\) and Taylor's formula give

\[
\boxed{t_A-t_0=\frac{\Delta_A(t_0,X)}{B_0}
 +O\left(\frac{X\Lambda^4D^2}{m^3}
                    +\frac{H_AD}{m^2}\right).}
\tag{NC16}\label{eq:ansatz-nc16}
\]

To check the error, insert \(h_{0}=\Delta(t_{0})/B_{0}\) into \(b_{lin}(t_{0}+h_{0})-\Delta(t_{0}+h_{0}).\) Its residual is bounded by \(C X \Lambda^{4} h_{0}^{2}+H_{A}|h_{0}|;\) division by the lower derivative bound proves NC16. For unscaled one-arm rows the nonlinear error is \(O(X^{-4}).\) Substituting NC6 into NC16 gives the actual finite-height leading corrections

\[
\boxed{t_{A_0}-t_0=-\frac{7\pi t_0}{16XB_0}
       +O(1/(X^3\lambda)+X^{-4}),}
\]
\[
\boxed{t_A-t_0=\frac{\pi t_0^2L}{32XB_0}
          +O(\Lambda^2/X^3),\quad A=A_1,A_J,A_S,A_4,A_{exp}.}
\tag{NC17}\label{eq:ansatz-nc17}
\]

Keeping B0 exact avoids imposing its coarser relative 1/logX error on these finite-height corrections. The elementary center is later; the corrected centers are earlier. Any pair in the corrected row differs by \(O(\Lambda^{2}/X^{3}).\) In particular all obey

\[
\frac{t_{A_0}-t_0}{w_X}
 =\frac{7\pi t_0}{4X\lambda}[1+O(1/\lambda)],
\]
\[
\frac{t_A-t_0}{w_X}
 =-\frac{\pi t_0^2L}{8X\lambda}[1+O(1/\lambda)],
\qquad t_A-t_0=o(w_X).
\tag{NC18}\label{eq:ansatz-nc18}
\]

The leading asymptotic forms in NC18 are explanatory; NC17 is the more accurate paid correction. For the factor-two distinction NC16 instead gives

\[
t_{2A}-t_A=\frac{\log2}{B_0}
                +O(1/(X^2\lambda^2)),\qquad
\frac{t_{2A}-t_A}{w_X}
 =-\frac{4\log2}{\lambda}[1+O(1/\lambda)].
\tag{NC19}\label{eq:ansatz-nc19}
\]

It is also sub-width, but only inverse-logarithmically. This shows why a first curvature correction is useful for arm accuracy without being necessary merely to achieve \(o(w_{X})\) centering. A fixed multiplicative amplitude change and an \(O(\log X/X)\) logarithmic change have different center effects and must use the same sensitivity NC15.

All assertions have a finite sufficient onset: impose the compact branch/source gates of Notes19--20; near-one ratios \(<1/4; \lambda\) large enough that \(32/\lambda<t_{0}-64<128/\lambda<1; C_{M} \lambda/X<1/4;\) and the displayed derivative/remainder smallness gates. Each follows eventually from the proved inequalities. No numeric X0, certified center query or affordable root-finding bound is supplied.

\subsubsection{What the displacement changes in the inherited layer}
\label{subsec:nc-5}

Parent Note28 JT33 uses exactly t0. Re-expressing its times around a new center means \(t=t_{A}-w_{X} \tau=t_{0}-w_{X}(\tau-(t_{A}-t_{0})/w_{X}).\) For bounded \(\tau,\) NC18 changes its leading normalized density profile by \(O(1/X),\) and NC19 by \(O(1/\log X),\) since \(k_{h}\) is Lipschitz and \(\tau\) -> \(\max(\tau,0)\) is Lipschitz. Both lie within the inherited \(O(1/\log X)\) allowance. Thus that inherited resolved law can be restated around these centers with the same error, but is not sharpened.

There is also a direct paid local layer-change bound. For each actual holomorphic arm \(a_{A}\) define \(\omega_{A}=\partial_{X} \arg a_{A}(t,X),\) retaining the coherent near-one \(\log; \omega_{lin}=\Im \partial_{X} \log G.\) On \(N_{M}, \omega_{lin}\) is \(\lambda/4+O(1)\) and \(b_{lin}=O_{M}(\lambda).\) The full-buffer log ratios of every unscaled one-arm model have size \(C \Lambda^{2}/X.\) Cauchy on radius \(c/\Lambda\) gives

\[
|\omega_A-\omega_{lin}|\le C\Lambda^3/X,\quad
|\Delta_A|\le C\Lambda/X.
\tag{NC20}\label{eq:ansatz-nc20}
\]

No derivative of a real modulus allowance is used. With \(d_{A}=\max(b_{A},0)/\omega_{A},\) and \(L_{A}=\omega_{A}[k_{h}(y-d_{A})+k_{h}(y+d_{A})]/(2\pi),\) the Lipschitz bounds and \(\omega_{A}\) comparable to \(\lambda\) give

\[
|d_A-d_{lin}|\le C_M\{ |\Delta_A|/\lambda
                         +|\omega_A-\omega_{lin}|/\lambda\},
\]
\[
|L_A-L_{lin}|\le C_M(|\Delta_A|+|\omega_A-\omega_{lin}|)
                    \le C_M\Lambda^3/X=o(1).
\tag{NC21}\label{eq:ansatz-nc21}
\]

Consequently the inherited \(|R_{H}-L_{lin}|\le C\) implies \(|R_{H}-L_{A}|\le C+o(1)\) on this actual switch neighborhood. Multiplying the modulus by two adds a bounded layer change instead, again leaving the \(O(1)\) rate. These are comparisons of layer expressions, not a new density theorem. The comparison centers do not redefine the historical \(t_{c},lin\) in JT26.

The first unpaid improved consumer is a vanishing absolute bound for \(R_{H}-L_{A}\) (or a specified next layer profile) with the fixed horizontal average, corrected phase, finite-period smoothing and adaptive terms retained. Parent JT24 discarded order-one frozen-profile/unit/rounding contributions. Small center displacement and \(X^{-7}\) field error do not pay that discarded contribution. A different useful consumer is the actual filtered density/current/time/covariance of the fixed P4 field, with no replacement by this two-layer shape. That consumer is proved in Theorem~\ref{thm:fixed-p4-consumer} below. No direct density-time, current, covariance, forcing, regulator or spatial-join rate is inferred by differentiating the inherited \(O(1).\)


\subsection[The complete fixed-P4 tuple]{The complete fixed-$P_4$ tuple}
\label{sec:fixed-p4-proof}
\begin{theorem}[Fixed-$P_4$ physical comparison on the compact switch]
\label{thm:fixed-p4-consumer}

Fix \(\ell,h,\eta>0.\) There exist finite C,X0, depending on these data and the named compact source/sector bounds, such that all statements below hold for \(63\le t\le 65, X=|x|\ge X_{0}\) and physical \(|y|\le 1.\) Use the complete averaging support, collars and rectangles \(r=\ell+2, b=1+h+2,\) with one further analytic unit. No numerical value of X0 is asserted. Let \(\Lambda=1+\log X, e=X^{-7}\) and \(A=|S(t,s)|+|S(t,1-s)|.\) Keep exactly

\[
\begin{gathered}
H_t=-H_{zz},\quad H(0,z)=\xi((1+iz)/2)/8,\quad s=(1+iz)/2,\\
P=P_4=\sum_{w=s,1-s}S(t,w)T(t,w),\quad T=1+C_2+C_4,
\end{gathered}
\tag{PT1}\label{eq:ansatz-pt1}
\]

The full canonical \(\theta\) source, Note19's fixed coefficient order \(J=2,\) chosen implicit gamma saddle and determinant branch are unchanged. The physical regulator is exactly \(q=\epsilon_{g}^{2}\) of parent Note20; the fixed spatial filter is \(K=k_{\ell} k_{h},\) with no time average or mass normalization. Write \(\varrho=\log q\) to distinguish it from the chart radius. For \(F=H,P\) define actual quantities

\[
D_F=|F|^2+q,\ S_F=\log D_F,\ w_F=q/D_F,\ \chi_F=w_F(1-w_F),
\quad \rho_F=\frac{q|F_z|^2}{\pi D_F^2},\quad
j_F=\frac{qF_{zz}\overline{F_z}}{\pi D_F^2},
\]
\[
U_F=K*S_F,\ W_F=K*w_F,\ R_F=K*\rho_F,\ J_F=K*j_F,
\quad Q_F=K*(S_{F,x}^2-S_{F,y}^2)-(U_{F,x}^2-U_{F,y}^2),
\]
\[
C_F=\Delta S_F/(4\pi)-\rho_F,\quad
b_{reg}=\varrho_t+(\varrho_{xx}+\varrho_x^2)/2,
\quad \mathcal S_F=(1-2w_F)(\varrho_t\rho_F+
                     \nabla\varrho\cdot(\Re j_F,\Im j_F)).
\tag{PT2}\label{eq:ansatz-pt2}
\]

Here \(C_{F}\) is the actual regulator curvature \(C_{reg}\) for a holomorphic F. Write \(\delta\) for H minus P in any observable. Every spatial derivative below is a real multi-index, at fixed physical time and coordinates.

\begingroup\small
\begin{longtable}{@{}>{\raggedright\arraybackslash}p{0.25\textwidth} >{\raggedright\arraybackslash}p{0.22\textwidth} >{\raggedright\arraybackslash}p{0.43\textwidth}@{}}
\toprule
Outputs & Spatial orders & Absolute bound\\
\midrule
\endhead
$\delta U,\delta W$ & $\le4$ & $Ce$\\
$\delta R,K*\delta C$ & $\le2$ & $Ce$\\
$\delta J,\delta Q$, including $\Delta\delta Q$ & $\le2$ & $Ce\Lambda^2$\\
$\delta J$ internally for source balance & $\le3$ & $Ce\Lambda^2$\\
$\delta U_t,\delta W_t,\delta R_t,K*\delta C_t$, $K*(b_{\rm reg}\delta w)$ and $K*\delta\mathcal S$ & $\le2$ & $CeX\Lambda^2$\\
$2\Re K*Z_0$ & $\le4$ & $Ce$\\
$K*\mathcal D_E$ & $\le2$ & $Ce$\\
\bottomrule
\end{longtable}
\endgroup

$Z_0$ and \(\mathcal D_E\) are the actual retained model residual corrections in PT9 and PT21 below. The $\Delta\delta Q$ assertion is the scalar Laplacian supplied by the spatial-order-two $Q$ bound; it is not differentiated twice again. The last two rows pay the heat-potential correction, its Laplacian through order two, and the density/current balance correction. All compatibility and field-join corrections listed in Section~\ref{subsec:pt-7} vanish identically for this unjoined holomorphic model.

\end{theorem}
\subsubsection{The exact compact source input}\label{subsec:pt-1}
The actual compact input is Ansatz Note20 RB5, not \(t=64\) inserted into an eventual onset:

\[
|\partial_z^j(H-P)|\le C A e\Lambda^j\ (0\le j\le3),\quad
|(H-P)_t|\le C A e\Lambda^2,\quad |E|\le C A e,
\]
\[
E=P_t+P_{zz}=-\sum_{w=s,1-s}S(t,w)\mathscr Q C_4(t,w).
\tag{PT3}\label{eq:ansatz-pt3}
\]

The exact operator $\mathscr Q$ and finite residual identity are Note19 FO23/FO30. All ordinary arithmetic, middle, connectors and infinite ends were paid before the compact field bound. The proof retains PT3's separate field, direct-time and residual inputs. No propagation from a small residual, assumed initial zero reality or fitted data is used.

\subsubsection{Acquire corrected phases, every collar and analytic unit}
\label{subsec:pt-2}

Work first near positive real height. Put $a_\pm(t,z)=S(t,(1\pm iz)/2)T(t,(1\pm iz)/2)$. Note19's finite FO8 formula, on the actual compact interval, gives

\[
C_2=O(X^{-3}),\ C_4=O(X^{-5}),\quad
T_w=O(X^{-4}),\ T_{ww}=O(X^{-5}).
\tag{PT4}\label{eq:ansatz-pt4}
\]

Indeed each monomial is a fixed power of t/B times products of \(\alpha^{(j)}(p)=O(X^{-j}), p_{w}=B^{-1}=O(1);\) differentiating once raises the inverse-height order by one. The sectorial derivatives through the finite order needed here are already available in Note19 at \(J=2.\) Thus \(|T-1|<1/4\) after an onset, so both corrected arms are zero-free on all the larger buffers and conjugate on the real axis.

Choose their coherent logarithms and define holomorphic functions

\[
\mathfrak A=\exp[(\log a_++\log a_-)/2],\quad
\Theta=(\log a_+-\log a_-)/(2i),\quad P=2\mathfrak A\cos\Theta.
\tag{PT5}\label{eq:ansatz-pt5}
\]

$\mathfrak A$ is positive and \(\Theta\) real on the real axis. Using \((\log S)_{w}=\alpha(p)+t \alpha''(p)/(4B^{2}),\) PT4 and the compact saddle expansion from Notes20/22 gives, on the full buffers,

\[
\Theta_z=\tfrac14\log(X/(4\pi))+O(\Lambda/X),\quad
\Theta_{zz}=O(X^{-1}),\quad
\mathfrak A_z/\mathfrak A=O(1),\qquad \omega=\Theta_x\asymp\Lambda.
\tag{PT6}\label{eq:ansatz-pt6}
\]

For example with \(D_{4}(w)=(\log S)_{w}+T_{w}/T, \Theta_{z}=[D_{4}(s)+D_{4}(1-s)]/4;\) the sum gives log-height growth while the difference defining the amplitude derivative is bounded. These are derivatives of exact holomorphic arms. The classical analytic inputs are \href{https://dlmf.nist.gov/1.9.E31}{Cauchy1.9.31}, \href{https://dlmf.nist.gov/1.10#iv}{Rouche and maximum modulus} and the \href{https://dlmf.nist.gov/5.11#ii}{sectorial gamma bounds}, rechecked in their primary DLMF statements on 26 October 2026. The bound $\Re\Theta_z\ge c\Lambda$ holds throughout the fixed rectangle for large X. Since \(\partial_{y} \Im\Theta=\Re\Theta_{z}\) and $\Im\Theta(x)=0$, a zero of $\cos\Theta$ there must be real. Real phase levels \((j+1/2)\pi\) have one simple zero \(c_{j}\) each and separation comparable to $1/\omega$. Their number intersecting any fixed horizontal support is \(O(\omega).\)

Choose nested disjoint phase disks of fixed sufficiently small radii below half the phase spacing. Their physical radii are comparable to \(1/\omega\) and fit inside the reserved analytic unit, including every disk meeting the filter support. On their complement, \(|\cos\Theta|\ge c \cosh(\Im\Theta).\) Since the corrected carrier sum is \(|a_+|+|a_-|=2|\mathfrak A|\cosh(\Im\Theta),\) this is a full transverse carrier-sum floor, not a real-axis amplitude bound. Corrected and uncorrected saddle envelopes are comparable by PT4.

Rouche on the large disks and PT3 exclude any additional H zeros in the complement and give exactly one H zero \(a_{j}\) per P disk. Real type forces that zero to be real. On the smaller radius \(C e/\omega\) circle, the model slope dominates PT3, giving

\[
|a_j-c_j|\le Ce/\omega,\quad
|F_z(\text{its root})|\ge c\mathfrak A(t,c_j)\omega,
\quad F=H,P.
\tag{PT7}\label{eq:ansatz-pt7}
\]

These conclusions acquire the P4 collars anew; they are not imported from the leading saddle's zero set. On nested scaled disks \(u=\omega_{c}(z-c_{j}),\) normalized by \(A_{c}=\mathfrak A(t,c_{j}),\) write \(P/A_{c}=u B(u), H/A_{c}=(u-e_{0})C(u), |e_{0}|\le Ce.\) On the larger boundary

\begin{quote}\small
\(C-B=(H-P)/[A_{c}(u-e_{0})]+e_{0}(P/A_{c})/[u(u-e_{0})]=O(e).\)\par
\end{quote}
Both units extend holomorphically through their respective roots. The maximum principle gives \(|C-B|\le Ce\) throughout the disk; the explicit cosine and compact bounds keep B and C above and below fixed positive moduli. Hence \(\log(C/B)=O(e)\) with the needed smaller-disk derivatives from Cauchy.

At each real root set $v_F=F_{zz}/F_z$. For $P$,
\[
 v_P=2\mathfrak A_z/\mathfrak A+\Theta_{zz}/\Theta_z=O(1).
\]
The spatial jets through three, root displacement and slope bounds in
PT3/PT7 give $|v_H-v_P|\le Ce\omega$, so both residues are bounded.
On the larger collars the exact pole decompositions are

\[
F_{zz}/F=v_F/(z-\text{root}_F)+g_F(z),\quad
|g_F|\le C\omega^2,\quad |g_H-g_P|\le Ce\omega^2.
\tag{PT8}\label{eq:ansatz-pt8}
\]

On the boundary the relative floor and jet bounds prove the difference bound. Subtract the two explicit poles, whose boundary difference has the same bound, then apply the maximum principle to the analytic remainder. This proves PT8 inside. Actual local root tracks, when needed, obey \(a_{t}=v_{H}\) and \(c_{t}=v_{P}-E(t,c)/P_{z}(t,c);\) their difference is \(O(e \omega+e/\omega).\) Algebraic residues have not been confused with model velocities. No population cutoff or collar indicator is differentiated.

\subsubsection{Prove the weak primitive and pole bounds for every positive q}
\label{subsec:pt-3}

Define smooth quantities, including their values at zeros, by

\[
V_{m,F}=w_F^m F_{zz}\overline F/D_F,
\qquad Z_m=w_P^m E\overline P/D_P,\qquad m=0,1,2.
\tag{PT9}\label{eq:ansatz-pt9}
\]

All L1 norms in this argument use planar area on a full fixed support buffer. We prove

\[
\boxed{\|\delta S\|_1+\|\delta w\|_1\le Ce,\quad
\|\delta V_m\|_1\le Ce\Lambda^2,\quad \|Z_m\|_1\le Ce.}
\tag{PT10}\label{eq:ansatz-pt10}
\]

The scalar argument holds for arbitrary positive spatial \(Q(u)=q/A_{c}^{2},\) held fixed pointwise while changing the field. For scalar f,g, away from their removable zero sets,

\begin{quote}\small
\(|\log[(|f|^{2}+Q)/(|g|^{2}+Q)]|\le 2|\log|f/g||,\)\par
\(|Q/(|f|^{2}+Q)-Q/(|g|^{2}+Q)|\le |\log|f/g||/2.\)\par
\end{quote}
Use the units above and integrate the translation of \(\log|u|\) along the root segment. Its gradient is $1/|u|$, uniformly integrable on the fixed disk, so the scaled log/weight integral costs Ce, not e log e. Restore physical area, sum \(O(\omega)\) disks and use the relative floor on the complement to obtain the first bound in PT10. The common physical \(\log|A_{c}|^{2}\) cancels only in the difference; no potential renormalization is performed. Polynomial weights in w, in particular \(\chi=w(1-w),\) inherit the same Ce comparison.

For the poles put

\begin{quote}\small
\(\Pi_{m}[f]=(u-a)^{-1}[Q/(|f|^{2}+Q)]^{m} |f|^{2}/(|f|^{2}+Q).\)\par
\end{quote}
Remove disks of radius C e around both roots. Their pole integrals are \(O(e).\) On the complement interpolate \(f_{s}=(u-s e_0)B \exp[s \log(C/B)].\) The bounded scalar weight and its derivative with respect to \(\log|f_{s}|\) give \(|\partial_{s} \Pi_{m}|\le Ce(|u-s e_0|^{-2}+|u-s e_0|^{-1}).\) Thus

\begin{quote}\small
\(\int |\Pi_{m}[(u-e_{0})C]-\Pi_{m}[uB]| \le Ce(1+|\log e|).\)\par
\end{quote}
No derivative or lower bound of Q enters. Physical area and the pole factor give per-collar translation \(Ce(1+|\log e|)/\omega.\) Summing collars costs \(Ce(1+|\log e|)=O(e \Lambda),\) since \(e=X^{-7}.\) Residue differences cost \(Ce \omega\) globally; PT8's regular parts and the changed scalar weights cost at most \(Ce \omega^{2}.\) The exterior is pointwise \(O(e \omega^{2}).\) This proves the second PT10 bound with the requested \(\Lambda^{2},\) including its logarithmic payment explicitly.

For the residual, E/P has residue \(E(t,c)/P_{z}(t,c)=O(e/\omega).\) On collar boundaries \(|E/P|\le Ce\) from the separate PT3 residual and floor. Subtract the pole and apply the maximum principle: its regular part is \(O(e).\) Multiplication by \(w^{m}(1-w)\) removes the singularity in the smooth \(Z_{m}\) circuit. The pole costs \(O(e/\omega^{2})\) per physical disk and the regular part the same; sum \(O(\omega)\) disks and the \(O(e)\) exterior to prove \(||Z_{m}||_1\le Ce.\) No \(E_{z}, E_{t}\) or inverse regulator estimate is needed in this weak bound.

\subsubsection{Pay the exact flat regulator before multiplying the weak bounds}
\label{subsec:pt-4}

On [63,65] the original regulator can be written exactly

\[
\varrho(t,u)=a(t)-\beta(t)u^2,\quad a=2\log E_0,\quad
\beta=0\ (t\le64),\quad \beta=\chi_{reg}/\kappa\ (t>64).
\tag{PT11}\label{eq:ansatz-pt11}
\]

The source saddle \(M_{0},\kappa\) and E0 are those of parent Notes19--20, not the gamma factor M or its complex saddle. Both a and \(\beta\) are smooth with bounded required derivatives on this compact interval. For \(\delta=t-64>0,\) Note22's exact flat-switch calculation gives \(\log \beta=-64/\delta+f(\delta),\) with f smooth at0. Direct differentiation therefore proves

\begin{quote}\small
\(|\beta'|\le C \beta[1+\log_+^{2}(1/\beta)]   (\beta>0),\)\par
\(\beta'=0 when \beta=0.\)\par
\end{quote}
Choose a fixed sufficiently large B0, for example \(B_{0}=\pi\) after increasing the height onset. In the regime \(\beta X\le B_{0},\) monotonicity of \(v(1+\log^{2}(1/v))\) for small positive v yields, on the full buffer,

\[
|\nabla\varrho|+|\Delta\varrho|\le C,\quad
|\varrho_t|+|b_{reg}|\le C X\Lambda^2,\quad
|\nabla\varrho_t|\le C\Lambda^2,\quad
|\Delta\varrho_t|\le C\Lambda^2/X.
\tag{PT12}\label{eq:ansatz-pt12}
\]

The exact derivatives are
\[
\begin{gathered}
\varrho_x=-2\beta u,\quad \varrho_{xx}=-2\beta,\quad
\varrho_t=a'-\beta'u^2,\quad \varrho_{xt}=-2\beta'u,\quad
\varrho_{xxt}=-2\beta',\\
q_t=q\varrho_t,\quad q_x=q\varrho_x,\quad
q_{xx}=q(\varrho_{xx}+\varrho_x^2),\quad
q_{xt}=q(\varrho_{xt}+\varrho_x\varrho_t).
\end{gathered}
\]
These are retained in the actual time circuits. Below $64$, $\beta=0$
and the derivatives of $a$ remain bounded; those times are included directly.

In the regime \(\beta X>B_{0},\) compact arm asymptotics and PT4 imply \(\log|a_+(t,u)|=-\pi u/8+O(\Lambda^{2}),\) uniformly for bounded \(|u-X|.\) Since \((X-r-1)^{2}\ge X^{2}/2\) eventually,

\[
\lambda_q:=\frac{\sup_{|u-X|\le r+1}\epsilon_g(t,u)}
                    {\mathfrak A(t,X)}\le Ce^{-cX},\qquad
|\log\lambda_q|\le C(1+X^2).
\tag{PT13}\label{eq:ansatz-pt13}
\]

Indeed \(\log\lambda_{q}\le -\beta X^{2}/4+\pi X/8+C \Lambda^{2},\) and \(B_{0}=\pi\) leaves a strictly negative linear term. All \(\varrho\) derivatives and the finite products used below grow at most polynomially in X.

On each collar the unit floor gives \(w_{F}\le a_{q}^{2}/(|z-root|^{2}+a_{q}^{2}), a_{q}=C \lambda_{q}/\omega;\) outside, \(w_{F}\le C \lambda_{q}^{2}.\) Radial integration and bounded residues give

\[
\int w_F\le C\lambda_q^2[1+(1+|\log\lambda_q|)/\omega],
\]
\[
\int|w_F F_{zz}/F|\le C\{\omega^2\lambda_q^2
               [1+(1+|\log\lambda_q|)/\omega]+\lambda_q\}.
\tag{PT14}\label{eq:ansatz-pt14}
\]

For \(w_{P} E/P\) replace \(\omega^{2}\) by e in the regular part and the pole coefficient by \(e/\omega.\) Thus every fixed polynomial coefficient times w, V1,V2 or $Z_1$,$Z_2$ is exponentially paid. V0 and $Z_0$ are not declared exponentially small: their PT10 comparison or residual bound is used where their coefficients are constant. Increase X0 so all required fixed polynomial times \(\exp(-cX)\) are below e. No global lower floor of q is introduced.

\subsubsection{Static, direct-time and curvature rows}
\label{subsec:pt-5}

Let \(||\psi||_{(k)}\) be the sum of test sup norms through order k. The exact weak identity for a holomorphic F is

\[
4\pi\langle\rho_F,\psi\rangle
=\int S_F\Delta\psi+2\int w_F\nabla\varrho\cdot\nabla\psi
 +\int[w_F\Delta\varrho+\chi_F|\nabla\varrho|^2]\psi.
\tag{PT15}\label{eq:ansatz-pt15}
\]

Subtract before taking absolute values. PT10 and PT12 give $|\langle\delta\rho,\psi\rangle|\le Ce\|\psi\|_{(2)}$. In the resolved regime the S difference is still Ce, while each coefficient-weighted \(w/\chi\) term is separately paid by PT14. Hence the same bound holds there. Using \(\psi=\partial^{\gamma} K, |\gamma|\le 2,\) proves the static R row. The static U,W rows follow directly from their L1 differences and bounded weak kernel derivatives through four.

For any F with \(F_{t}=-F_{zz}+E_{F},\) direct differentiation gives

\[
S_{F,t}=-2\Re V_{0,F}+\varrho_t w_F+2\Re Z_{0,F},
\]
\[
w_{F,t}=\varrho_t\chi_F+2\Re V_{1,F}-2\Re Z_{1,F},
\]
\[
\chi_{F,t}=\varrho_t\chi_F(1-2w_F)
  +2\Re(V_{1,F}-2V_{2,F})-2\Re(Z_{1,F}-2Z_{2,F}).
\tag{PT16}\label{eq:ansatz-pt16}
\]

Here \(Z_{m,H}=0\) and \(Z_{m,P}=Z_{m}\) of PT9. These identities use the actual field equations; they are not obtained by differentiating PT10. PT10/PT12, or the individual weighted bounds PT14, prove

\[
\|\delta S_t\|_1+\|\delta w_t\|_1+
       \|\delta\chi_t\|_1\le C\tau_X,\qquad
\tau_X=e X\Lambda^2.
\tag{PT17}\label{eq:ansatz-pt17}
\]

For example \(\varrho_{t} \delta w\) costs exactly \(Ce X \Lambda^{2}\) in the unresolved regime. In the resolved regime \(\varrho_{t} w\) and all weighted V1/$Z_1$ terms are polynomial times exponentially small; the unweighted \(\delta V_{0}\) and $Z_0$ remain \(O(e \Lambda^{2})\) and \(O(e).\) Convolution proves the requested \(U_{t},W_{t}\) orders<=2. PT17 also pays their spatial derivatives internally through four if needed, by moving derivatives to K; no extra time derivative is taken.

Differentiate PT15 at fixed physical coordinates and fixed test:

\[
\begin{split}
4\pi\langle\rho_{F,t},\psi\rangle={}&\int S_{F,t}\Delta\psi
 +2\int(w_{F,t}\nabla\varrho+w_F\nabla\varrho_t)\cdot\nabla\psi\\
 &+\int[w_{F,t}\Delta\varrho+w_F\Delta\varrho_t
 +\chi_{F,t}|\nabla\varrho|^2
 +2\chi_F\nabla\varrho\cdot\nabla\varrho_t]\psi.
\end{split}
\tag{PT18}\label{eq:ansatz-pt18}
\]

At each finite point \(q>0\) makes this an ordinary differentiation of smooth circuits on compact support. Uniformity is then supplied by PT12/PT14/PT17. Every subtracted term is \(\le C \tau_{X}||\psi||_{(2)}.\) This is a direct density-time proof, independent of current balance. With \(|\gamma|\le 2,\) the total kernel order is \(\le 4,\) proving \(R_{t}.\)

For curvature, directly subtract PT15 from $\langle\Delta S,\psi\rangle$:
\[
\begin{split}
4\pi\langle C_F,\psi\rangle={}&-2\int w_F\nabla\varrho\cdot\nabla\psi\\
 &-\int\bigl[w_F\Delta\varrho+\chi_F|\nabla\varrho|^2\bigr]\psi.
\end{split}
\]
Its direct time derivative has the last two lines of PT18 with the minus sign and without the \(S_{t}\) term. It proves static Ce and time \(C \tau_{X}\) curvature comparisons through spatial order two, using at most three kernel derivatives. In particular both \(C_{t}\) and $\square C$ remain in the final equation. Lastly \(b_{reg} \delta w,\) with the same two-regime payments, proves its requested \(C \tau_{X}\) filtered row.

\subsubsection{Current, full covariance and actual signed source}
\label{subsec:pt-6}

The exact current identity, proved by differentiating the smooth fraction \(F_{zz} conjugate(F)/D_{F}\) even across zeros, is

\[
\pi j_F=\partial_{\bar z}V_{0,F}
                      +(\partial_{\bar z}\varrho)V_{1,F}.
\tag{PT19}\label{eq:ansatz-pt19}
\]

In a test pairing the first derivative transfers to \(\psi. PT10/PT12\) or PT14 therefore give a weak order-one current difference \(Ce \Lambda^{2}||\psi||_{(1)}.\) Consequently \(\partial^{\gamma} \delta J\) has that rate for all \(|\gamma|\le 3;\) internal order3 costs K order4. No division by \(F_{z}\) between roots occurs.

For \(G_{F}=S_{F,x}^{2}-S_{F,y}^{2}\) retain the exact signed spatial identity

\[
G_F=4\Re V_{0,F}-\square S_F
                        +(\varrho_{xx}+\varrho_x^2)w_F,
\quad \square=\partial_x^2-\partial_y^2.
\tag{PT20}\label{eq:ansatz-pt20}
\]

It follows from \(\square|F|^{2}=4\Re(F_{zz} \overline F)\) and \(\square \log D+S_{x}^{2}-S_{y}^{2}=(\square D)/D.\) Subtract PT20 before bounding: PT10 and the regulator split prove \(\partial^{\gamma} K*\delta G =O(e \Lambda^{2})\) through \(|\gamma|\le 2,\) using K through order4. Separately bounding the two positive gradient squares is not used.

To pay the actual mean products, write \(\nabla S_{F}=(2\Re L_{F},-2\Im L_{F})+w_{F} \nabla\varrho, L_{F}=(1-w_{F})F_{z}/F.\) Exterior floors give \(|F_{z}/F|\le C \omega;\) on collars \(F_{z}/F=1/(z-root)+O(\omega).\) The simple poles have uniformly integrable planar mass, with \(O(\omega)\) disks of radius \(O(1/\omega);\) their sum is \(O(1).\) The exterior and analytic units cost \(O(\omega).\) The regulator term is bounded by PT12 or PT14. Thus \(\|\nabla S_F\|_1\le C \Lambda\) and \(|\partial^{\nu} U_{F}|\le C \Lambda\) for \(1\le |\nu|\le 3.\) Static mean differences are Ce. The full Leibniz sum therefore gives \(\partial^{\gamma} \delta(U_{x}^{2}-U_{y}^{2})=O(e \Lambda), |\gamma|\le 2.\) Subtract these actual full means from \(K*G\) to obtain \(\delta Q=O(e \Lambda^{2}),\) including both second derivatives and \(\Delta Q.\) The physical additive potential is retained; covariance is not defined from an equation or replaced by component means.

For P let its actual residual density correction be

\[
\mathcal D_E=\frac{2q\Re(E_z\overline{P_z})}{\pi D_P^2}
                   -4\rho_P\Re Z_0.
\tag{PT21}\label{eq:ansatz-pt21}
\]

Vary PT15 in direction E at fixed q. Since $\delta S=2\Re Z_0$, $\delta w=-2\Re Z_1$ and $\delta\chi=-2\Re(Z_1-2Z_2)$, this proves exactly

\[
\begin{split}
4\pi\langle\mathcal D_E,\psi\rangle={}&2\Re\int Z_0\Delta\psi
 -4\Re\int Z_1\nabla\varrho\cdot\nabla\psi\\
 &-2\Re\int[Z_1\Delta\varrho
                  +(Z_1-2Z_2)|\nabla\varrho|^2]\psi.
\end{split}
\tag{PT22}\label{eq:ansatz-pt22}
\]

PT10/PT12 or PT14 gives $|\langle\mathcal D_E,\psi\rangle|\le Ce\|\psi\|_{(2)}$. The raw \(E_{z}\) term in PT21 is paid by this equivalent signed weak expression, not deleted. Thus \(K*\mathcal D_E\) through order2 and \(2\Re K*Z_{0}\) through order4 have the requested Ce bounds.

Finally direct differentiation of the actual density/current circuits gives
$\rho_{F,t}+\operatorname{div}j_F=\mathcal S_F+\mathcal D_{E_F}$,
with $\mathcal D_{E_H}=0$. Therefore, with $\delta=H-P$,
\[
 K*\delta\mathcal S=\delta R_t+\operatorname{div}\delta J+K*\mathcal D_E.
\]
PT18 is already proved independently. With internal current order3 and PT22, spatial orders<=2 of this actual signed source difference are \(C \tau_{X}.\) No integral of the source's absolute raw integrand and no product of filtered means is substituted for the signed source.

\subsubsection{Full inventory, equations and maximum kernel order}
\label{subsec:pt-7}

For this single model the compatibility field $\partial_{\bar z}P=0$. Thus the general GS8/GS12 quantities \(T_{C},(T_{C})_{t},T_{J},\Xi,D_{B}\) are identically zero, as are their requested spatial derivatives. \(C_{tot}=C_{reg},\) whereas \(E=E_{4}\) is retained and generally nonzero: $T_U=2\Re Z_0$ and \(\mathcal D_E\) are exactly the quantities paid above. Spatial, moving and temporal field-join addends vanish because there is no field join, time splice or moving filter in this packet.

The model equations, with actual Q and the unchanged q/K, are

\[
U_t=-\tfrac12(\square U+U_x^2-U_y^2+Q)
                +K*(b_{reg}w)+2\Re K*Z_0,\quad
\Delta U=4\pi(R+K*C),
\]
\[
\begin{split}
R_t={}&-\tfrac12\square R-\tfrac1{8\pi}\Delta(U_x^2-U_y^2+Q)
 +\tfrac1{4\pi}\Delta K*(b_{reg}w)\\
 &-K*C_t-\tfrac12\square K*C+\tfrac1{2\pi}\Delta\Re K*Z_0,
\end{split}
\]
\[
R_t+\operatorname{div}J=K*\mathcal S+K*\mathcal D_E.
\tag{PT23}\label{eq:ansatz-pt23}
\]

H has the same equations with residual corrections zero. Every component used in these displayed equations and in Note59's fixed full-tuple inventory is now paid at the stated orders. The internal V0 and \(c_{r} w\) rows through spatial order two follow directly from PT10 and PT12/PT14, where \(c_{r}=\varrho_{xx}+\varrho_{x}^{2}.\) Fixed-time and fixed observation-boundary traces inherit the displayed uniform filtered bounds; no moving-boundary transport theorem is asserted. The density equation is asserted at its displayed scalar order; differentiating the whole equation twice would demand extra covariance derivatives not requested here. The separately proved \(R_{t}\) row through order2 uses PT18 instead and does not incur those extra derivatives.

\begingroup\small
\begin{longtable}{@{}>{\raggedright\arraybackslash}p{0.23\textwidth} >{\raggedright\arraybackslash}p{0.24\textwidth} >{\raggedright\arraybackslash}p{0.24\textwidth} >{\raggedright\arraybackslash}p{0.18\textwidth}@{}}
\toprule
Consumer & Largest test differentiation & Largest filtered order & Maximum kernel order\\
\midrule
\endhead
U,W; heat-potential correction & 0 & 4 & 4\\
R, direct \(R_{t},\) residual \(\mathcal D_E\) & 2 & 2 & 4\\
\(C,C_{t}\) from their own weak identities & 1 & 2 & 3\\
J, including source's internal current & 1 & 3 & 4\\
signed covariance stress & 2 & 2 & 4\\
actual mean products & gradients transferred once, then derivatives & 3 on U & 3\\
\bottomrule
\end{longtable}
\endgroup

The fixed \(k(v)=(315/256)(1-v^{2})^{4}\) on \(|v|\le 1\) is C3 with a bounded piecewise-polynomial weak fourth derivative and no boundary \(\delta\) through order4. Products and scalings give the required K bounds. Pairings with its weak derivatives follow by smoothing or weak integration by parts. No fifth kernel derivative is used. Apply the exact identities to the whole translated test before splitting estimates into collars/exterior and regulator regimes. The splits are inequalities, not differentiated test partitions; there are no omitted collar-boundary, moving or temporal commutators.

\subsubsection{Onset, limits and source screening}
\label{subsec:pt-8}

Choose X0 after fixing \(\ell,h,\eta.\) Impose the compact Notes19--20 branch/field gates, PT4 near-one gates, \(\Re\Theta_{z}\ge c \Lambda,\) the nested-collar and Rouche smallness conditions, \(e \omega<1, B_{0}/X\le 1/2\) (the monotonicity bound is used only where \(\beta X\le B_{0}),\) and the finite list of polynomial \(\exp(-cX)\le X^{-7}\) absorptions used above. All are eventual inequalities with fixed constants and therefore give a finite sufficient onset. The theorem claims no numeric height or useful cost at that height. Negative real charts follow by evenness and the even physical regulator.

The physical rates improve the inherited parent Note31 PE44 main comparison rows by five height powers at this fixed compact scope, and its \(X^{-3}\) residual corrections by four. The full exact-q tuple, actual covariance, both curvature derivatives and all nonzero model defects remain. No estimate in the frozen requested rows is unpaid. This neither sharpens Note28's simplified-layer \(O(1)\) absolute/ \(O(1/\log X)\) normalized law nor reassembles the all-time construction. Notes58--61 retain their separate source/evaluator/weak-filter height-cost obligations; a regulator-floor-free approximation proof is not a finite-node filter algorithm or a useful height-cost theorem. The bounded-height compact core is outside this high-height physical result;
the field atlas of Theorem~\ref{thm:positive-atlas} has its own separate scope.

The source-specific input is the complete canonical xi/gamma/ordinary- arithmetic comparison PT3. Euler-factor and finite-prime perturbations lack the unchanged full coefficients; weighted-square, finite-cutoff and atomic controls lack the full canonical completion; reflected, smooth and positive-primitive controls lack the same complete source and arithmetic identity. The weak pole and physical transfer lemmas are generic and may apply to simple collars of controls too. They are not claimed to reject every derived control bound or prove global noncollision/RH. This is the premise screen required by Control12.


The tuple theorem closes the fixed compact $P_4$ consumer, with all
weak derivatives, full mean products, regulator curvature and actual
nonzero residual corrections retained. It does not sharpen the older
simplified-layer law, improve the global spatial/time joins, construct
a complete weak-primitive filter evaluator, or acquire numerical
constants/onsets. The canonical source premise enters through the
complete field theorem; the generic collar/weak lemmas exclude no
collision control by themselves. No RH, global noncollision or
simplicity theorem follows.
